\chapter{Combinational Logic Design}

\section{Continuous Assignments}
Combinational circuits produce outputs that depend only on the present values of their inputs. In
Verilog, the most direct way to describe this kind of hardware is by using a continuous assignment.
A continuous assignment represents a permanent connection from an expression to a net, similar to a
wire driven by a gate or a small network of gates in a schematic.\\

\noindent The keyword used for this purpose is \inlinecode{assign}. It is written inside a module,
but outside procedural blocks such as \inlinecode{initial} and \inlinecode{always}. The left-hand
side of a continuous assignment is normally a net such as \inlinecode{wire}, and the right-hand side
is an expression constructed from input signals, constants, and Verilog operators. Whenever any
signal on the right-hand side changes, the expression is re-evaluated and the driven output is
updated by the simulator. This continuous evaluation behavior matches the physical behavior of
combinational logic, where gate outputs respond to changes at their inputs after propagation delay.\\

\noindent The simplest example is a two-input AND gate:

\begin{codeblock}{verilog}
module and_gate(A, B, Y);
    input A, B;
    output Y;

    assign Y = A & B;
endmodule
\end{codeblock}

\noindent The output \inlinecode{Y} is driven by the bitwise AND of \inlinecode{A} and
\inlinecode{B}. If either input changes during simulation, \inlinecode{Y} is automatically updated.
No explicit event control or procedural statement is required because the assignment is concurrent
with the rest of the module.\\

\noindent An OR gate follows the same pattern:

\begin{codeblock}{verilog}
module or_gate(A, B, Y);
    input A, B;
    output Y;

    assign Y = A | B;
endmodule
\end{codeblock}

\noindent In this case, the operator \inlinecode{|} maps to a bitwise OR function. For one-bit
signals it behaves as the OR gate used in a logic diagram. For vectors, the same operator applies
bit by bit across corresponding positions of the operands.\\

\noindent Continuous assignments may also contain compound expressions. The following module drives
three outputs from the same set of inputs:

\begin{codeblock}{verilog}
module simple_logic(A, B, C, Y1, Y2, Y3);
    input A, B, C;
    output Y1, Y2, Y3;

    assign Y1 = A & B;
    assign Y2 = A | C;
    assign Y3 = (A & B) | (~C);
endmodule
\end{codeblock}

\noindent Each \inlinecode{assign} statement describes a separate concurrent hardware connection.
The expression for \inlinecode{Y3} contains an AND operation, an inversion, and an OR operation; a
synthesis tool interprets this as the corresponding combinational network. The order of the
statements in the source does not imply execution order. All continuous assignments are active at
the same time, which is why they are commonly used for RTL descriptions of combinational logic.

\section{Dataflow Modeling}
Dataflow modeling describes a circuit by specifying how data values are transformed as they move
from inputs to outputs. Instead of instantiating individual gates, the designer writes Boolean or
arithmetic expressions that define the required output behavior. The synthesis tool then constructs
the equivalent logic hardware from those expressions.\\

\noindent This is different from a structural description. In structural modeling, the source code
explicitly connects gates or lower-level modules, much like a textual schematic. In dataflow
modeling, the source code focuses on the relationship between signals. For example, a designer may
write the expression for an output directly rather than listing every AND, OR, and NOT gate needed
to realize it. This makes dataflow modeling concise and well suited to RTL design, especially for
combinational blocks whose behavior is naturally expressed by equations.\\

\noindent A simple dataflow example is the implementation of an inverter and a two-input gate:

\begin{codeblock}{verilog}
module basic_dataflow(A, B, Y_not, Y_and);
    input A, B;
    output Y_not, Y_and;

    assign Y_not = ~A;
    assign Y_and = A & B;
endmodule
\end{codeblock}

\noindent The unary operator \inlinecode{~} describes bitwise inversion, while \inlinecode{&}
describes bitwise AND. Since all signals are one bit wide in this example, each operator maps
directly to the corresponding logic gate.\\

\noindent More useful combinational functions are usually written as expressions that combine
several operators:

\begin{codeblock}{verilog}
module combinational_network(A, B, C, D, F);
    input A, B, C, D;
    output F;

    assign F = (A & B) | (C & ~D);
endmodule
\end{codeblock}

\noindent The output \inlinecode{F} becomes high when \inlinecode{A} and \inlinecode{B} are both
high, or when \inlinecode{C} is high while \inlinecode{D} is low. Parentheses are used to make the
intended grouping clear. Although Verilog has defined operator precedence rules, explicit grouping
is preferred in laboratory code because it reduces ambiguity during review.\\

\noindent Dataflow descriptions also work naturally with buses. In the next example, the same
operation is applied across four bits:

\begin{codeblock}{verilog}
module vector_logic(A, B, mask, Y);
    input [3:0] A, B, mask;
    output [3:0] Y;

    assign Y = (A & mask) | (B & ~mask);
endmodule
\end{codeblock}

\noindent Here, \inlinecode{A}, \inlinecode{B}, \inlinecode{mask}, and \inlinecode{Y} are four-bit
vectors. The bitwise operators act independently on each bit position. This style is common in RTL
because it describes a repeated logic pattern without writing a separate equation for every bit.\\

\noindent The following table summarizes the operators most frequently used in dataflow models of
combinational logic.

\begin{table}[H]
\centering
\renewcommand{\arraystretch}{1.2}
\begin{tabular}{|c|c|p{7cm}|}
\hline
\textbf{Operator} & \textbf{Name} & \textbf{Typical Hardware Interpretation} \\
\hline
\inlinecode{~} & Bitwise NOT & Inverter on each bit of the operand \\
\hline
\inlinecode{&} & Bitwise AND & AND gate for each bit position \\
\hline
\inlinecode{|} & Bitwise OR & OR gate for each bit position \\
\hline
\inlinecode{^} & Bitwise XOR & XOR gate for each bit position \\
\hline
\inlinecode{+} & Addition & Combinational adder network \\
\hline
\inlinecode{?:} & Conditional & Multiplexer selection logic \\
\hline
\end{tabular}
\caption{Common Verilog operators used in dataflow descriptions.}
\label{tab:dataflow_operators}
\end{table}

\noindent Dataflow modeling is widely used because it is compact, synthesizable, and close to the
Boolean equations used during circuit design. Its main limitation is that very complex expressions
can become difficult to read and debug. For larger designs, dataflow statements are often combined
with hierarchy: small combinational functions are written as equations, and larger systems are built
by instantiating those verified modules.

\section{Conditional Operators}
The conditional operator provides a compact way to select one of two expressions based on a
condition. It is also called the ternary operator because it uses three operands: the condition, the
expression selected when the condition is true, and the expression selected when the condition is
false. Its general form is:

\begin{codeblock}{verilog}
condition ? expression1 : expression2;
\end{codeblock}

\noindent If \inlinecode{condition} evaluates to true, the result of the whole expression is
\inlinecode{expression1}. If the condition evaluates to false, the result is
\inlinecode{expression2}. In synthesizable combinational RTL, this operator is normally interpreted
as selection logic. For one-bit selection, that hardware is a multiplexer.\\

\noindent A simple output selection example is shown below:

\begin{codeblock}{verilog}
module conditional_example(A, B, sel, Y);
    input A, B, sel;
    output Y;

    assign Y = sel ? A : B;
endmodule
\end{codeblock}

\noindent When \inlinecode{sel} is high, \inlinecode{Y} follows \inlinecode{A}. When
\inlinecode{sel} is low, \inlinecode{Y} follows \inlinecode{B}. The statement is continuously
evaluated, so any change on \inlinecode{sel}, \inlinecode{A}, or \inlinecode{B} is reflected at the
output.\\

\noindent The selected expressions may also be logic expressions rather than simple signals:

\begin{codeblock}{verilog}
module conditional_logic(A, B, C, mode, Y);
    input A, B, C, mode;
    output Y;

    assign Y = mode ? (A & B) : (A | C);
endmodule
\end{codeblock}

\noindent In this example, \inlinecode{mode} selects between two different logic functions. When
\inlinecode{mode} is high, the output is the AND of \inlinecode{A} and \inlinecode{B}. When
\inlinecode{mode} is low, the output is the OR of \inlinecode{A} and \inlinecode{C}. This is useful
when a circuit must route one of several possible results to a common output.\\

\noindent Conditional operators are also useful for signal routing on buses:

\begin{codeblock}{verilog}
module bus_select(data_a, data_b, sel, data_out);
    input [3:0] data_a, data_b;
    input sel;
    output [3:0] data_out;

    assign data_out = sel ? data_a : data_b;
endmodule
\end{codeblock}

\noindent The selection is applied to the complete four-bit vector. Hardware synthesis produces four
parallel one-bit selection circuits controlled by the same \inlinecode{sel} signal. This pattern is
the natural Verilog form of a multiplexer, which is developed in detail in the next section.

\section{Multiplexers}
A multiplexer, commonly abbreviated as MUX, is a combinational circuit that selects one input from
several available data inputs and forwards the selected value to a single output. The selection is
controlled by one or more select lines. Multiplexers are fundamental building blocks in digital
systems because they implement controlled routing: one data source is chosen while the others are
ignored.\\

\noindent A multiplexer has three main parts: data inputs, select inputs, and an output. The number
of select lines determines how many data inputs can be addressed. With one select line, two inputs
can be selected. With two select lines, four inputs can be selected. In general, \inlinecode{n}
select lines can address \inlinecode{2^n} input choices. Since the output is determined only by the
current values of the data and select inputs, a multiplexer is a combinational circuit and is well
suited to continuous assignments and dataflow modeling.

\subsection{2:1 Multiplexer}
A 2:1 multiplexer has two data inputs, one select input, and one output. When the select input is
low, the output follows \inlinecode{I0}. When the select input is high, the output follows
\inlinecode{I1}. The complete truth table is shown in Table~\ref{tab:mux2_truth}. The selected
input determines the output value, while the unselected input has no effect on the output.\\

\noindent The Boolean expression for a 2:1 multiplexer is:
\[
Y = \overline{sel}\cdot I0 + sel\cdot I1
\]

\noindent The first product term enables \inlinecode{I0} when \inlinecode{sel = 0}, and the second
product term enables \inlinecode{I1} when \inlinecode{sel = 1}.

\begin{table}[H]
\centering
\renewcommand{\arraystretch}{1.2}
\begin{tabular}{|c|c|c|c|}
\hline
\textbf{sel} & \textbf{I0} & \textbf{I1} & \textbf{Y} \\
\hline
0 & 0 & 0 & 0 \\
\hline
0 & 0 & 1 & 0 \\
\hline
0 & 1 & 0 & 1 \\
\hline
0 & 1 & 1 & 1 \\
\hline
1 & 0 & 0 & 0 \\
\hline
1 & 0 & 1 & 1 \\
\hline
1 & 1 & 0 & 0 \\
\hline
1 & 1 & 1 & 1 \\
\hline
\end{tabular}
\caption{Full truth table of a 2:1 multiplexer.}
\label{tab:mux2_truth}
\end{table}

\noindent The conditional operator gives the clearest Verilog description of this selection:

\begin{codeblock}{verilog}
module mux2_conditional(I0, I1, sel, Y);
    input I0, I1, sel;
    output Y;

    assign Y = sel ? I1 : I0;
endmodule
\end{codeblock}

\noindent The expression reads directly from the truth table. If \inlinecode{sel} is high,
\inlinecode{I1} is selected; otherwise \inlinecode{I0} is selected. This is a dataflow description
because the output is defined by an expression rather than by explicit gate instances.\\

\noindent The same 2:1 multiplexer can also be written using its Boolean equation:

\begin{codeblock}{verilog}
module mux2_dataflow(I0, I1, sel, Y);
    input I0, I1, sel;
    output Y;

    assign Y = (~sel & I0) | (sel & I1);
endmodule
\end{codeblock}

\noindent This version exposes the underlying AND-OR structure. The term \inlinecode{~sel & I0}
enables \inlinecode{I0} when the select signal is low, and the term \inlinecode{sel & I1} enables
\inlinecode{I1} when the select signal is high. The OR operation combines the enabled path into the
output.

\subsection{4:1 Multiplexer}
A 4:1 multiplexer extends the same idea to four data inputs. It requires two select lines because
four choices must be encoded. The select input is usually represented as a two-bit vector
\inlinecode{sel[1:0]}, where \inlinecode{sel[1]} is the most significant select bit and
\inlinecode{sel[0]} is the least significant select bit. The output follows exactly one of the four
data inputs for each select combination.\\

\noindent The Boolean expression for a 4:1 multiplexer is:
\[
Y = \overline{sel[1]}\cdot\overline{sel[0]}\cdot I0
  + \overline{sel[1]}\cdot sel[0]\cdot I1
  + sel[1]\cdot\overline{sel[0]}\cdot I2
  + sel[1]\cdot sel[0]\cdot I3
\]

\noindent Each product term corresponds to one decoded select condition. Only the selected data
input is allowed to contribute to \inlinecode{Y}; the other data inputs are masked by the select
logic.

\begin{table}[H]
\centering
\renewcommand{\arraystretch}{1.2}
\begin{tabular}{|c|c|c|c|c|c|c|}
\hline
\textbf{sel[1]} & \textbf{sel[0]} & \textbf{I0} & \textbf{I1} & \textbf{I2} & \textbf{I3} & \textbf{Y} \\
\hline
0 & 0 & 0 & X & X & X & 0 \\
\hline
0 & 0 & 1 & X & X & X & 1 \\
\hline
0 & 1 & X & 0 & X & X & 0 \\
\hline
0 & 1 & X & 1 & X & X & 1 \\
\hline
1 & 0 & X & X & 0 & X & 0 \\
\hline
1 & 0 & X & X & 1 & X & 1 \\
\hline
1 & 1 & X & X & X & 0 & 0 \\
\hline
1 & 1 & X & X & X & 1 & 1 \\
\hline
\end{tabular}
\caption{Full truth table of a 4:1 multiplexer.}
\label{tab:mux4_truth}
\end{table}

\noindent A compact dataflow implementation can be written by nesting conditional operators:

\begin{codeblock}{verilog}
module mux4_conditional(I0, I1, I2, I3, sel, Y);
    input I0, I1, I2, I3;
    input [1:0] sel;
    output Y;

    assign Y = sel[1] ? (sel[0] ? I3 : I2) : (sel[0] ? I1 : I0);
endmodule
\end{codeblock}

\noindent The outer conditional checks \inlinecode{sel[1]}. If \inlinecode{sel[1]} is low, the
selection is between \inlinecode{I0} and \inlinecode{I1}; if \inlinecode{sel[1]} is high, the
selection is between \inlinecode{I2} and \inlinecode{I3}. The inner conditional in each branch then
uses \inlinecode{sel[0]} to choose the final input. This corresponds to a tree of 2:1 multiplexers.\\

\noindent The same function can be expressed using minterms:

\begin{codeblock}{verilog}
module mux4_dataflow(I0, I1, I2, I3, sel, Y);
    input I0, I1, I2, I3;
    input [1:0] sel;
    output Y;

    assign Y = (~sel[1] & ~sel[0] & I0) |
               (~sel[1] &  sel[0] & I1) |
               ( sel[1] & ~sel[0] & I2) |
               ( sel[1] &  sel[0] & I3);
endmodule
\end{codeblock}

\noindent This form is longer, but it makes the Boolean decoding of the select lines explicit. Each
product term corresponds to one row of the truth table, and only one data input is enabled for each
binary value of \inlinecode{sel}.\\

\noindent Multiplexers are used extensively in digital systems. They select operands for arithmetic
units, choose between register outputs, route data onto internal buses, select next-state values in
control logic, and implement configurable datapaths. The conditional operator and dataflow
assignments introduced earlier provide the most direct Verilog notation for these functions.

\section{Full Adder Design}
A Half Adder adds two one-bit inputs and produces a Sum and Carry output. It is useful for
introducing combinational design, but it cannot accept a carry from a lower-order bit position.
Multi-bit addition requires that each bit position include the carry generated by the previous
position. This is the purpose of a Full Adder.\\

\noindent A Full Adder has three one-bit inputs: \inlinecode{A}, \inlinecode{B}, and
\inlinecode{Cin}. The input \inlinecode{Cin} represents the carry-in from a previous stage. It has
two outputs: \inlinecode{Sum}, which is the one-bit addition result, and \inlinecode{Cout}, which is
the carry-out passed to the next stage. The complete behavior is shown in
Table~\ref{tab:full_adder_truth}.

\begin{table}[H]
\centering
\renewcommand{\arraystretch}{1.2}
\begin{tabular}{|c|c|c|c|c|}
\hline
\textbf{A} & \textbf{B} & \textbf{Cin} & \textbf{Sum} & \textbf{Cout} \\
\hline
0 & 0 & 0 & 0 & 0 \\
\hline
0 & 0 & 1 & 1 & 0 \\
\hline
0 & 1 & 0 & 1 & 0 \\
\hline
0 & 1 & 1 & 0 & 1 \\
\hline
1 & 0 & 0 & 1 & 0 \\
\hline
1 & 0 & 1 & 0 & 1 \\
\hline
1 & 1 & 0 & 0 & 1 \\
\hline
1 & 1 & 1 & 1 & 1 \\
\hline
\end{tabular}
\caption{Truth table of the Full Adder.}
\label{tab:full_adder_truth}
\end{table}

\noindent From the truth table, the Sum output is high when an odd number of inputs are high. The
Carry output is high when at least two of the three inputs are high. The Boolean equations are:
\[
\text{Sum} = A \oplus B \oplus Cin
\]
\[
\text{Cout} = (A \cdot B) + (B \cdot Cin) + (A \cdot Cin)
\]

\noindent The carry equation shows why the Full Adder is suitable for multi-bit arithmetic. A carry
is produced either because \inlinecode{A} and \inlinecode{B} are both high, or because one operand
bit is high while the incoming carry is also high. In a ripple carry adder, this
\inlinecode{Cout} becomes the \inlinecode{Cin} of the next more significant Full Adder stage. This
stage-to-stage dependency is called carry propagation.\\


\noindent A Full Adder can also be understood as two Half Adders and an OR gate. The first Half Adder
adds \inlinecode{A} and \inlinecode{B}, producing an intermediate sum and carry. The second Half
Adder adds the intermediate sum to \inlinecode{Cin}. The two carry outputs are then ORed to produce
\inlinecode{Cout}. In RTL, the same behavior can be written more directly using continuous
assignments:

\begin{codeblock}{verilog}
module full_adder(A, B, Cin, Sum, Cout);
    input A, B, Cin;
    output Sum, Cout;

    assign Sum = A ^ B ^ Cin;
    assign Cout = (A & B) | (B & Cin) | (A & Cin);
endmodule
\end{codeblock}

\noindent The module declaration defines the Full Adder as a reusable design block. The input
declaration lists the three one-bit operands, including the carry-in. The output declaration lists
the two result signals. The first continuous assignment implements the three-input XOR function for
\inlinecode{Sum}. The second continuous assignment implements the carry equation using AND terms
for the possible carry-generating input pairs and an OR operation to combine them. Finally,
\inlinecode{endmodule} closes the module so that it can be compiled independently or instantiated
inside a larger design.\\

\noindent An alternative compact implementation uses Verilog addition and concatenation:

\begin{codeblock}{verilog}
module full_adder_compact(A, B, Cin, Sum, Cout);
    input A, B, Cin;
    output Sum, Cout;

    assign {Cout, Sum} = A + B + Cin;
endmodule
\end{codeblock}

\noindent The concatenation \inlinecode{{Cout, Sum}} forms a two-bit result. Since the maximum value
of \inlinecode{A + B + Cin} is three, two bits are sufficient to represent both the carry and the
sum. This form is concise and synthesizable, but the explicit Boolean version is often preferred
when first studying the circuit because it shows the logic structure directly.\\

\noindent The Full Adder is the basic cell used to construct wider adders. In the following
laboratory exercise, several Full Adder stages will be connected so that the carry-out of each stage
feeds the carry-in of the next stage, forming a 4-bit ripple carry adder.

\section{Laboratory Exercise -- 4-Bit Ripple Carry Adder}

\subsection{Objective}
The objective of this laboratory exercise is to design, simulate, and verify a 4-bit Ripple Carry
Adder using Verilog HDL. The exercise combines the combinational design techniques introduced in
this chapter and applies them to a practical arithmetic circuit.\\

\noindent By completing this exercise, the reader will demonstrate hierarchical design, module
instantiation, carry propagation between adder stages, and functional verification through
simulation. The design will be written as reusable Verilog modules and verified in the Cadence
simulation environment using a separate testbench.

\subsection{Theory}
Binary addition is performed bit by bit, beginning with the least significant bit and progressing
toward the most significant bit. At each bit position, the circuit must add the corresponding bits
of the two operands and also include any carry generated by the previous lower-order position. A
Half Adder is insufficient for this task because it has no carry-in input.\\

\noindent A Full Adder solves this limitation by accepting three inputs: \inlinecode{A},
\inlinecode{B}, and \inlinecode{Cin}. It produces a \inlinecode{Sum} bit and a carry-out signal
\inlinecode{Cout}. When several Full Adders are connected together, each stage handles one bit of a
multi-bit addition. The carry-out of one stage becomes the carry-in of the next stage. This
carry-chain structure is called a ripple carry adder because the carry information ripples from the
least significant stage to the most significant stage.\\

\noindent The ripple carry adder is simple and well suited for laboratory study. It clearly shows
module reuse, signal interconnection, and the relationship between a one-bit arithmetic cell and a
wider arithmetic circuit.

\subsection{Ripple Carry Adder Architecture}
A 4-bit Ripple Carry Adder is constructed from four one-bit Full Adder modules connected in series.
The first Full Adder adds \inlinecode{A[0]}, \inlinecode{B[0]}, and the external
\inlinecode{Cin}. Its carry-out is connected to the carry-in of the second Full Adder. This pattern
continues until the fourth stage produces the final carry-out \inlinecode{Cout}.\\

\noindent The input operands are \inlinecode{A[3:0]} and \inlinecode{B[3:0]}. The least significant
bits are \inlinecode{A[0]} and \inlinecode{B[0]}, while the most significant bits are
\inlinecode{A[3]} and \inlinecode{B[3]}. The output \inlinecode{Sum[3:0]} contains the four result
bits, and \inlinecode{Cout} indicates whether the addition generated a carry beyond the most
significant bit.

% TODO: Insert Ripple Carry Adder block diagram
\begin{figure}[H]
    \centering
    \includegraphics[width=0.6\textwidth]{resources/ripple_carry.png}
    \caption{Conceptual block diagram of a 4-bit Ripple Carry Adder constructed from four Full Adder stages.}
    \label{fig:ripple_carry_adder_block}
\end{figure}

\noindent Data enters all four stages through the operand buses, but the carry path is sequential
through the combinational chain. The first stage computes \inlinecode{Sum[0]} and an internal carry
\inlinecode{c1}. The second stage uses \inlinecode{c1} to compute \inlinecode{Sum[1]} and
\inlinecode{c2}. The third stage produces \inlinecode{Sum[2]} and \inlinecode{c3}. The fourth stage
uses \inlinecode{c3} to compute \inlinecode{Sum[3]} and the final \inlinecode{Cout}. Although there
is no clock in this circuit, the output cannot settle fully until the carry has propagated through
the required stages.

\subsection{Full Adder Module}
The Full Adder module used in this exercise is the same one-bit combinational block discussed
earlier in the chapter. It is written using continuous assignments so that the Boolean equations
are represented directly in Verilog.

\begin{codeblock}{verilog}
module full_adder(A, B, Cin, Sum, Cout);
    input A, B, Cin;
    output Sum, Cout;

    assign Sum = A ^ B ^ Cin;
    assign Cout = (A & B) | (B & Cin) | (A & Cin);
endmodule
\end{codeblock}

\noindent The inputs \inlinecode{A} and \inlinecode{B} are the operand bits for the current stage,
and \inlinecode{Cin} is the carry entering that stage. The output \inlinecode{Sum} is generated by
the XOR of the three inputs. The output \inlinecode{Cout} is generated when any two or more inputs
are high, as expressed by the three AND terms combined by OR operations. This module is designed as
a reusable cell and will be instantiated four times in the 4-bit adder.

\subsection{4-Bit Ripple Carry Adder Implementation}
The top-level Ripple Carry Adder module demonstrates hierarchical combinational design. Instead of
rewriting the Full Adder equations four times, the design instantiates four copies of the
\inlinecode{full_adder} module and connects them through intermediate carry wires.

\begin{codeblock}{verilog}
module ripple_carry_adder_4bit(A, B, Cin, Sum, Cout);
    input [3:0] A, B;
    input Cin;
    output [3:0] Sum;
    output Cout;

    wire c1, c2, c3;

    full_adder FA0 (
        .A(A[0]),
        .B(B[0]),
        .Cin(Cin),
        .Sum(Sum[0]),
        .Cout(c1)
    );

    full_adder FA1 (
        .A(A[1]),
        .B(B[1]),
        .Cin(c1),
        .Sum(Sum[1]),
        .Cout(c2)
    );

    full_adder FA2 (
        .A(A[2]),
        .B(B[2]),
        .Cin(c2),
        .Sum(Sum[2]),
        .Cout(c3)
    );

    full_adder FA3 (
        .A(A[3]),
        .B(B[3]),
        .Cin(c3),
        .Sum(Sum[3]),
        .Cout(Cout)
    );
endmodule
\end{codeblock}

\noindent The module interface declares two four-bit input operands, one external carry-in, a
four-bit sum output, and one final carry-out. The internal wires \inlinecode{c1}, \inlinecode{c2},
and \inlinecode{c3} carry information between adjacent Full Adder stages.\\

\noindent Instance \inlinecode{FA0} is the least significant stage. It adds \inlinecode{A[0]},
\inlinecode{B[0]}, and \inlinecode{Cin}, then produces \inlinecode{Sum[0]} and internal carry
\inlinecode{c1}. Instance \inlinecode{FA1} uses \inlinecode{c1} as its carry-in and produces
\inlinecode{c2}. Instance \inlinecode{FA2} uses \inlinecode{c2} and produces \inlinecode{c3}.
Finally, instance \inlinecode{FA3} uses \inlinecode{c3} and produces the most significant sum bit
and final \inlinecode{Cout}. This carry chaining is the defining feature of the ripple carry
architecture.

\subsection{Testbench Development}
A testbench is required because the Ripple Carry Adder module does not generate its own input
stimulus. The testbench instantiates the design under test, applies representative operand values,
waits for the combinational outputs to settle, and records the results for inspection in the
simulator and waveform viewer.\\

\noindent The following testbench checks additions without carry generation, additions that produce
internal carries, and additions that overflow beyond four bits.

\begin{codeblock}{verilog}
`timescale 1ns/1ps

module ripple_carry_adder_4bit_tb;
    reg [3:0] A, B;
    reg Cin;
    wire [3:0] Sum;
    wire Cout;

    ripple_carry_adder_4bit dut (
        .A(A),
        .B(B),
        .Cin(Cin),
        .Sum(Sum),
        .Cout(Cout)
    );

    initial begin
        $monitor($time, " A=%b B=%b Cin=%b Cout=%b Sum=%b",
                 A, B, Cin, Cout, Sum);

        A = 4'b0000; B = 4'b0000; Cin = 1'b0; #10;
        A = 4'b0011; B = 4'b0100; Cin = 1'b0; #10;
        A = 4'b0111; B = 4'b0001; Cin = 1'b0; #10;
        A = 4'b0101; B = 4'b0011; Cin = 1'b1; #10;
        A = 4'b1111; B = 4'b0001; Cin = 1'b0; #10;
        A = 4'b1111; B = 4'b1111; Cin = 1'b1; #10;

        $finish;
    end
endmodule
\end{codeblock}

\noindent The first test applies zero operands and verifies the reset-like baseline behavior of the
combinational circuit. The second test performs addition without overflow. The third and fourth
tests require carry propagation through one or more stages. The fifth test produces a carry beyond
the four-bit sum field, and the final test verifies the maximum input condition with an asserted
carry-in. The \inlinecode{#10} delays separate each stimulus vector in simulation time so that the
waveform viewer clearly shows the response to each input combination. The \inlinecode{$finish}
system task terminates the simulation after all vectors have been applied.

\subsection{Simulation Procedure}
The simulation should be performed as a controlled laboratory workflow. Keep the Full Adder source,
Ripple Carry Adder source, and testbench source in the same project directory or add all three files
explicitly to the Cadence simulation setup.

\begin{enumerate}
    \item Create a new Cadence project or working directory for the Ripple Carry Adder exercise.
    \item Create the Verilog source file containing the \inlinecode{full_adder} module.
    \item Create the Verilog source file containing the \inlinecode{ripple_carry_adder_4bit} module.
    \item Create the testbench file containing \inlinecode{ripple_carry_adder_4bit_tb}.
    \item Add all design and testbench files to the simulation environment.
    \item Compile the source files and confirm that there are no syntax or elaboration errors.
    \item Launch the simulation with \inlinecode{ripple_carry_adder_4bit_tb} as the top-level module.
    \item Add \inlinecode{A}, \inlinecode{B}, \inlinecode{Cin}, \inlinecode{Sum}, and
    \inlinecode{Cout} to the waveform viewer.
    \item Run the simulation long enough for all test vectors to execute.
    \item Compare the waveform and console output against the expected binary addition results.
\end{enumerate}

% TODO: Insert source compilation screenshot
% TODO: Insert simulation launch screenshot
% TODO: Insert waveform viewer screenshot

\noindent During compilation, errors are most commonly caused by mismatched module names, missing
source files, or incorrect port names during instantiation. During simulation, incorrect results are
usually caused by a broken carry connection between stages or a mismatch between the expected and
actual bit ordering of the buses.

\subsection{Expected Results}
The simulation should show that \inlinecode{Sum[3:0]} contains the lower four bits of the binary
addition and that \inlinecode{Cout} contains the carry beyond the most significant bit. For example,
\inlinecode{0011 + 0100} should produce \inlinecode{Sum = 0111} and \inlinecode{Cout = 0}. The
case \inlinecode{1111 + 0001} should produce \inlinecode{Sum = 0000} and \inlinecode{Cout = 1},
showing overflow beyond the four-bit output range.\\

\noindent Carry propagation should be visible in the functional behavior of the result. When a
lower-order stage generates a carry, the next stage must include that carry in its own sum
calculation. Although the internal carry wires may not be top-level ports, they can be added to the
waveform viewer from the design hierarchy if the simulator setup allows internal signal probing.

% TODO: Insert Ripple Carry Adder waveform
\begin{figure}[H]
    \centering
    \fbox{\parbox{0.86\linewidth}{\centering TODO: Insert Ripple Carry Adder waveform}}
    \caption{Expected waveform for the 4-bit Ripple Carry Adder showing operand inputs, carry-in, sum output, and carry-out.}
    \label{fig:ripple_carry_adder_waveform}
\end{figure}

\noindent The waveform verifies correct operation when each applied input vector produces the
corresponding binary sum after the simulator advances past the stimulus delay. A mismatch in
\inlinecode{Sum} indicates an error in the Full Adder equation or bit connection. A mismatch in
\inlinecode{Cout} usually indicates an error in the final carry connection or in one of the
intermediate carry wires.

\subsection{Conclusion}
This exercise demonstrated the design and verification of a 4-bit Ripple Carry Adder using Verilog
HDL. The circuit was built hierarchically by defining a reusable Full Adder module and instantiating
it four times inside a top-level combinational design. The intermediate carry wires illustrated how
carry propagation links individual one-bit stages into a wider arithmetic circuit.\\

\noindent The testbench verified the design by applying representative input combinations and
checking the resulting sum and carry-out behavior through simulation. This workflow reinforces the
standard RTL design method used throughout digital systems engineering: describe a reusable module,
connect modules hierarchically, compile the design, simulate with controlled stimulus, and verify
the waveform against expected behavior. The same discipline will be used in later chapters when the
manual moves from combinational circuits to sequential logic and more complex digital systems.
